\section*{Chapter \ref{ch:pl:the-risk-grid} --- The Risk Grid}

\begin{solution}{exo:pl:the-risk-grid:1}
2\,200 tasks (220 trade batches times ten scenario batches), adding $2\,200 \times 5 = 11\,000$ core-seconds of fixed cost, about 10\% of the
pricing.
\end{solution}

\begin{solution}{exo:pl:the-risk-grid:2}
The finish is set by one task --- the slow trade's batch across all scenarios, about 5\,380~seconds --- and by when it starts. More cores
start it earlier, since the queue in front of it drains faster, but cannot shorten it: on 64 cores it still ends 6\,922~seconds after the
start.
\end{solution}

\begin{solution}{exo:pl:the-risk-grid:3}
Sixteen repricings (each of eight pillars up and down) plus the base: 17 calls; the adjoint records one valuation and runs one reverse sweep,
at a cost of a few valuations whatever the number of pillars.
\end{solution}

\begin{solution}{exo:pl:the-risk-grid:4}
The number of tasks grows as the batches shrink --- 44\,000 at five scenarios --- and each carries five seconds of fixed cost; below about 50
scenarios the fixed costs grow faster than the balance improves.
\end{solution}

\begin{solution}{exo:pl:the-risk-grid:5}
The failing swaption sits in a trade batch of several thousand trades; every task of that batch --- one per scenario batch --- fails, so the
batch's results under every scenario are lost: 38.3~million cells.
\end{solution}

\begin{solution}{exo:pl:the-risk-grid:6}
When the market-data version changes, since every result's key changes. Intraday risk on a new market is covered by sensitivities carried
forward (chapter~18) and, if needed, a partial revaluation of the trades most exposed to the move.
\end{solution}

\begin{solution}{exo:pl:the-risk-grid:7}
On 64 cores the naive plan meets the window from 5 to 500 scenarios a task (best 50: 2\,319~seconds) and fails only at 1\,000; the
cost-aware plan meets it at every granularity, 5 to 1\,000 (best 250: 1\,989~seconds). More cores widen the good range; they do not remove the need to plan.
\end{solution}

\begin{solution}{exo:pl:the-risk-grid:8}
Doubling the cores moved the finish from 06:51 to 06:25, still late, because the lateness came from one task and the plan, not from capacity.
Measure the tasks, isolate what is slow, choose the granularity; buy capacity only when the lower bound of a good plan is above the window.
\end{solution}

\begin{solution}{pb:pl:the-risk-grid:1}
\begin{enumerate}
\item 105\,000 trades by 1\,000 scenarios: 105~million pricing calls (plus the base valuations).
\item The 5\,000 Monte Carlo autocallables: 22.2~milliseconds a call, 99\% of the 112.2~seconds of pricing per scenario.
\item Five seconds: loading the snapshot and the trades, writing the results.
\item 90~minutes, 04:30 to 06:00, on 32 cores.
\item About 4\,000~seconds at 100 scenarios a task: the work over 32 cores.
\item 2.35, 1.19 and 2.93~hours.
\item 20, 50, 100 and 250 scenarios a task (20 by sixteen seconds).
\item It makes one task of more than 5\,300~seconds, started late because its estimate is small.
\item It uses last night's measured cost to give the slow trade tasks of their own, fanned out finely.
\item 05:34, with 250 scenarios a task.
\item Per task: two retries, then the task's results are lost; per trade: the trade is quarantined with a report and the task completes.
\item 38.3~million results and 1\,100 core-seconds per task; 1\,000 results and nothing per trade.
\item 5.9 times (388 against 66~microseconds).
\item 1.02~million bumped repricings replaced by 60\,000 recorded valuations.
\item 1.5~million calls for 1\,500 changed trades, against 105~million.
\item \textbf{Named result.} On 32 cores the grid takes 2.35~hours with a task per trade batch across all scenarios, 1.19 with 100 scenarios
a task and 2.93 with five; only 20 to 250 scenarios meet the 90-minute window with the naive plan, and a cost-aware plan meets it at every
granularity from 20 to 1\,000 (best 1.07~hours at 250); doubling the cores moves the naive plan from 06:51 to 06:25; adjoints replace 1.02~million
bumped repricings of the swaps by 60\,000 valuations, and an incremental intraday rerun makes 1.5~million calls instead of 105~million.
\item Plan by measured cost, with the scenario batch chosen from the fixed cost and the cluster, and quarantine failing trades.
\item Each task's duration against its estimate, the slowest trades, failures and quarantines, the finish against the window, and the lower
bound.
\item When a good plan's lower bound, not its \omterm{def:pl:distributed-compute-and-schedulers:makespan}{makespan}, approaches the window.
\item Its longest tasks and when they start.
\end{enumerate}
\end{solution}

\begin{solution}{iq:pl:the-risk-grid:1}
The tasks on the critical path: which ran last and longest, and whether their durations matched their estimates; then failures and retries;
then the granularity against the fixed cost of a task; capacity last.
\iqlookfor{Critical path before capacity.}
\end{solution}

\begin{solution}{iq:pl:the-risk-grid:2}
Batch trades by measured cost to a target per scenario, batch scenarios so that tasks last minutes (many times the fixed cost, much less than
the window), isolate and fan out expensive trades, and schedule longest first.
\iqlookfor{Cost-based batching, two dimensions.}
\end{solution}

\begin{solution}{iq:pl:the-risk-grid:3}
Because it gives every sensitivity of a price for a few times the cost of the price, where bumping costs two prices per sensitivity; it does
not help when there are few sensitivities, when the code cannot be taped (a closed-source pricer), or for non-smooth payoffs without smoothing.
\iqlookfor{Cost independent of the number of inputs.}
\end{solution}

\begin{solution}{iq:pl:the-risk-grid:4}
Quarantine it: the pricing loop catches its failure, records it in an error report with its reason, completes the task without it, and the
report goes to the owner of the trade's data before the next run.
\iqlookfor{Isolate the failure, never the task.}
\end{solution}

\begin{solution}{iq:pl:the-risk-grid:5}
A grid planned from measured costs on a scheduler with fair share and longest-first ordering; tasks sized from their fixed costs; results
stored by trade, trade version and market version; failures quarantined with reports; adjoint sensitivities where the library supports them;
incremental intraday reruns for changed trades and sensitivities carried forward for market moves.
\iqlookfor{Planning, robustness, reuse.}
\end{solution}
